% General note: This is not the place to describe the changes in detail;
% that should be done either in the introduction or through the change log.
% Provide only the highest level summary here.

\ifdraftspec
\paragraph*{202X Draft Specification, XXXX, 202X.}
This document contains a draft of the \MPI/ specification as of the date of
publication.
It has not been adopted as an official \MPI/ specification, and
is provided for comment only.
This document includes updates and new features that will be present in the
final \mpivdoti/ document.
\else
\fi

\paragraph*{Version 5.0: June 5, 2025.}
This version of the \MPIV/ standard is a major update and includes
significant new functionality.
The largest change is the addition of a standard Application
Binary Interface (ABI) to allow interoperability of different
implementations.

\paragraph*{Version 4.1: November 2, 2023.}
This version of the \MPIIVDOTI/ standard contains mostly corrections
and clarifications to the \MPIIVDOTO/ document. Several routines, the
attribute key \mpiconstskip{MPI\_HOST}, and the \code{mpif.h} Fortran
include file are deprecated.

\paragraph*{Version 4.0: June 9, 2021.}
This version of the \MPIIV/ standard is a major update and includes
significant new functionality.
The largest changes are the addition of large-count versions of many
routines to address the limitations of using an \code{int}
or \code{INTEGER} for the count parameter, persistent collectives, partitioned
communications, an alternative way to initialize \MPI/, application
info assertions, and improvements to the definitions of error handling.
In addition, there are a number of smaller improvements and corrections.

\paragraph*{Version 3.1:  June 4, 2015.}
This document contains mostly corrections and clarifications to the
\mpiiiidoto/ document.  The largest change is a correction to the Fortran
bindings introduced in \mpiiiidoto/.
Additionally, new functions added
include routines to manipulate \code{MPI\_Aint} values in a portable manner, nonblocking
collective I/O routines, and routines to get the index value by name for
\mpiskipfunc{MPI\_T} performance and control variables.

% It would be a terrible mistake to list all changes to MPI-3 here.  If
% this list is unsatisfactory, the next best option is to list *no*
% specific enhancements, and instead simply instruct the reader that there
% are many extensions.
%
\paragraph*{Version 3.0:  September 21, 2012.}
Coincident with the development of \mpiiidotii/, the \MPI/ Forum began
discussions of a major extension to \MPI/.  This document contains the
\MPIIII/ standard.  This version of the \mpiiii/ standard
contains significant extensions to \MPI/ functionality, including
nonblocking collectives, new one-sided communication operations, and
Fortran 2008 bindings.  Unlike \mpiiidotii/, this standard is
considered a major update to the \MPI/ standard.  As with previous
versions, new features have been adopted only when there were
compelling needs for the users.  Some features, however, may have more
than a minor impact on existing \MPI/ implementations.

\paragraph*{Version 2.2:  September 4, 2009.}
This document contains mostly corrections and clarifications to the
\mpiiidoti/ document.  A few extensions have been added; however all
correct \mpiiidoti/ programs are correct
\mpiiidotii/ programs.  New features
were adopted only when there were compelling needs for users, open
source implementations, and minor impact on existing \MPI/
implementations.  

% the double 2008 is a typo, but not corrected in the moment, because the official standard should be printed 
% \paragraph*{Version 2.1: \today, 2008.}
\paragraph*{Version 2.1: June 23, 2008.}
This document combines the previous documents 
\mpiidotiii/ (May 30, 2008) 
and \mpiiidoto/ (July 18, 1997). 
Certain parts of \mpiiidoto/, such as some sections of Chapter~4, 
Miscellany, and Chapter~7, Extended Collective
Operations, have  
been merged into the chapters of \mpiidotiii/. Additional errata and 
clarifications collected by the \mpi/ Forum are also included in 
this document.

\paragraph*{Version 1.3: May 30, 2008.}
This document combines the previous 
documents \mpiidoti/ (June 12, 1995) and the \mpiidotii/ chapter in \mpiii/ 
(July 18, 1997). Additional errata collected by the \mpi/ Forum 
referring to \mpiidoti/ and \mpiidotii/ are also included in this 
document.

\paragraph*{Version 2.0: July 18, 1997.}
Beginning after the release of 
\mpiidoti/, the \mpi/ Forum began meeting to consider corrections and 
extensions. \mpiii/ has been focused on process creation and 
management, one-sided communications, extended collective 
communications, external interfaces and parallel I/O. A miscellany 
chapter discusses items that do not fit elsewhere, in particular 
language interoperability.

\paragraph*{Version 1.2: July 18, 1997.}
The \mpiii/ Forum introduced \mpiidotii/ as 
Chapter~3 
in the standard ``\mpiii/: Extensions to the Message-Passing
Interface'', July 18, 1997.  This section contains clarifications 
and minor corrections to Version 1.1 of the \mpi/ standard. The only 
new function in \mpiidotii/ is one for identifying to which version of 
the \mpi/ standard the implementation conforms. There are small 
differences between \mpii/ and \mpiidoti/. There are very few differences 
between \mpiidoti/ and \mpiidotii/, but large differences between \mpiidotii/
and \mpiii/.

\paragraph*{Version 1.1: June, 1995.}
Beginning in March, 1995, the Message-Passing Interface Forum reconvened to 
correct errors and make clarifications in the \mpi/ document of May 5, 1994,
referred to below as Version~1.0.  These discussions resulted in 
Version~1.1.  The changes from Version~1.0 are minor.  A version
of this document with all changes marked is available.  

%MPx-1.2
\paragraph*{Version 1.0:  May, 1994.}
The Message-Passing Interface Forum, with participation from over
40 organizations, has been meeting since January 1993 to discuss and
define a set of library interface standards for message
passing.
The Message-Passing Interface Forum
is not sanctioned or supported by any official
standards organization.

The goal of the Message-Passing Interface, simply stated, is to
develop a widely used
standard for writing message-passing programs.
As such the interface should
establish a practical, portable, efficient, and flexible standard 
for message-passing.

This is the final report, Version~1.0, of
the Message-Passing Interface Forum.
This document contains all the
technical features proposed for the interface.  This copy of the draft
was processed by \LaTeX{} on
May 5, 1994.

%% WDG - I'm ok with just deleting this, as it can be viewed as either attached 
%% to Version 1.0 (in which case it should be left) or as a separate point,
%% in which case, it should be deleted (but *not* replaced with a newer
%% version)
